\section*{Results}

\subsection*{A negative freshwater transport, and the two factors that set it}

Recomputed from monthly fields on the contiguous Atlantic section
(Methods), the record-mean $\Fovs$ at 34.5\degS{} is negative in all three
reanalyses (Fig.~\ref{fig:sign}a,b): $-0.038$~Sv in ORAS5 over 1958--2025,
$-0.084$~Sv in GLORYS12V1 over 1993--2025 and $-0.132$~Sv in ECCO-V4r4 over
1992--2017. The sign is not marginal. It holds in 61 of 68 annual means in
ORAS5, 32 of 33 in GLORYS12V1 and 26 of 26 in ECCO, and the 95\% circular
block-bootstrap interval excludes zero in every product at every block length
from 2 to 20 years, so the result does not depend on that choice. Because a
mean over a trending record is not an estimate of any well-defined population
quantity, we also give the most recent decade, which is the quantity a reader
interested in the present state should use: $-0.088$~[$-0.094$, $-0.080$]~Sv in
ORAS5 over 2016--2025, $-0.132$~[$-0.141$, $-0.121$] in GLORYS12V1 and
$-0.128$~[$-0.136$, $-0.119$] in ECCO over 2008--2017.

A fourth product, SODA 3.15.2, gives $-0.009$~Sv with an interval that includes
zero, but we do not count it either way. The archive we can obtain provides
annual rather than monthly fields, and $\Fovs$ is bilinear in velocity and
salinity, so forming it from annual means discards the within-year covariance.
Measured directly in ECCO, where both routes are available, that omission shifts
$\Fovs$ by $+0.012$~Sv, larger in magnitude than SODA's entire record mean. The
same test gives $+0.000$~Sv in ORAS5 and $-0.001$~Sv in GLORYS12V1, so the bias
is product-dependent and small in the two eddy-permitting products. SODA can
neither support nor refute the sign, and we say so rather than counting it.

The magnitudes span a factor of three across the three products, which is
larger than any within-product interval. The reanalyses differ more about how
much freshwater the overturning carries than about which way it carries it, and
only the second statement is robust.

Equation~\eqref{eq:factorisation} splits each of these numbers into a transport
and a salinity contrast. The northward limb carries 18.7~Sv in ORAS5, 21.3~Sv
in GLORYS12V1 and 14.6~Sv in ECCO, against limb salinity contrasts of $+0.068$,
$+0.136$ and $+0.307$~PSU. (The factorisation is exact year by year, so the
product of the record means differs slightly from the mean of the products, by
6\% in ORAS5.) Almost the entire three-product spread in $\Fovs$ is in
$\Delta S$, which varies by a factor of 4.5, not in $T$, which varies by a
factor of 1.5. Whatever the reanalyses constrain at this section, it is the
transport far better than the salinity structure.

\subsection*{The decline outruns the transport by an order of magnitude}

ORAS5 and GLORYS12V1 both decline over their own records, by $-1.37$~mSv per
year ($p<10^{-4}$, $N_{\text{eff}}=23.7$) and $-4.36$~mSv per year
($p=0.017$, $N_{\text{eff}}=7.2$). ECCO does not decline at all
($+0.11$~mSv per year, $p=0.76$, $N_{\text{eff}}=23.1$, minimum detectable
trend 0.75~mSv per year). Over the 1993--2017 window all three share, no
product has a significant \emph{$\Fovs$} trend, but that is a statement about
power rather than about the ocean: the smallest trend those 25 years could have
resolved is 1.3~mSv per year in ORAS5 and 7.1 in GLORYS12V1, and the GLORYS12V1
point estimate over the common window ($-4.76$) is larger in magnitude than over
its own record. The northward-limb transport behaves differently over that same
window, and in the opposite direction to the paper's framing: it declines in all
three products, by $-0.97$~Sv per decade in ORAS5 ($p=0.052$,
$N_{\text{eff}}=10.1$), $-0.75$ in GLORYS12V1 ($p=0.036$, $N_{\text{eff}}=14.3$)
and $-0.35$ in ECCO ($p=0.062$, $N_{\text{eff}}=21.3$). Over the one window in
which the three products are directly comparable it is the transport, not
$\Fovs$, that has the near-significant trends.

The cleanest statement of the result needs no ratio at all. Expressed as
fractional rates, $\Fovs$ is changing at $-36$\% of its own magnitude per decade
in ORAS5 and $-52$\% in GLORYS12V1, while the transport that carries it changes
at $-2.6$\% and $-0.6$\% per decade. Over the Argo era the same comparison gives
$-35$\% against $-0.7$\% and $-20$\% against $+0.8$\%. In those windows the two quantities are
changing at rates that differ by between one and two orders of magnitude, which
is the argument in its simplest form and is free of any denominator that might
approach zero. The separation is not uniform across every window: over the
common 1993--2017 window, where the $\Fovs$ trend is not resolvable and the
transport trend nearly is, the ORAS5 ratio narrows to 2.5. Across the four
windows we examine, the transport rate is between 1\% and 132\% of the $\Fovs$
rate in the two declining products. The wide end comes entirely from windows in
which the $\Fovs$ trend is not resolvable: 40\% from the common window in
ORAS5 and 132\% from the eleven-year clean window in GLORYS12V1, where the
$\Fovs$ trend has the wrong sign and $p=0.79$. Over the windows in which the
denominator is resolvable the range is 1\% to 7\%.

Splitting the trend by equation~\eqref{eq:factorisation} puts numbers on it
(Fig.~\ref{fig:attribution}b). In ORAS5 the total is $-1.37$~mSv per year, of
which the transport term contributes $+0.09$ [$+0.06$, $+0.13$], the salinity
term $-1.46$ [$-1.77$, $-1.14$] and a covariance residual $+0.002$, 3\% of the
transport term. In GLORYS12V1 the total is $-4.36$, the transport term $+0.05$
[$-0.12$, $+0.21$], the salinity term $-4.36$ [$-5.78$, $-2.94$] and the
residual $-0.05$, which is 108\% of the transport term and larger than it. The
transport share of the trend is 0.068 [0.046, 0.094] in ORAS5 and 0.012
[0.001, 0.051] in GLORYS12V1; assigning the residual entirely to the transport
instead of dropping it moves the GLORYS12V1 share to 0.001 and flips its sign, so for that product the honest statement is that the transport
contribution is not distinguishable from zero rather than that it is
positive. The transport term opposes the total with probability 1.00 in ORAS5
over its own record but only 0.73 in GLORYS12V1, and 0.36 in GLORYS12V1 over
the Argo era, where its point estimate ($-0.09$~mSv per year) has the same sign
as the total. The magnitude claim is robust; the sign claim is not universal
and we do not make it universal.

What the salinity term represents is unusually specific
(Fig.~\ref{fig:mechanism}a). The contrast $\Delta S$ widens because the
northward limb becomes saltier while the southward limb does not change.
$S_{\uparrow}$ rises by $+0.025$~PSU per decade in ORAS5 ($p<10^{-4}$,
$N_{\text{eff}}=24.2$) and $+0.071$ in GLORYS12V1 ($p=0.017$,
$N_{\text{eff}}=6.9$, which is under five degrees of freedom and should be read
accordingly), whereas $S_{\downarrow}$ changes by $-0.002$ ($p=0.06$) and
$-0.001$ ($p=0.24$) over the same records. The southward limb is not
uniformly flat: over the matched 1993--2017 window the ORAS5 value is
$-0.003$ and does reach significance ($p=0.006$). The asymmetry rests on
magnitude, $S_{\uparrow}$ moving six times faster, rather than on
$S_{\downarrow}$ being exactly stationary. Over the Argo era the two products converge, to $+0.051$
($p=0.009$) and $+0.035$ ($p=0.005$); note that this is a rise for ORAS5 and a
halving for GLORYS12V1, so the agreement improves while the GLORYS12V1 rate
falls. Over the matched 1993--2017 window neither is individually significant
($+0.019$, $p=0.11$, $N_{\text{eff}}=15.8$; $+0.083$, $p=0.10$,
$N_{\text{eff}}=5.2$) and they differ by a factor of four.

The depth-resolved salinity trends behind these numbers disagree far more than
the limb-weighted quantities do (Fig.~\ref{fig:mechanism}b). Through the upper
500~m GLORYS12V1 trends at $+0.10$ to $+0.13$~PSU per decade against a
layer-mean $+0.009$ in ORAS5, which crosses zero near 200~m, so the two differ
by more than an order of magnitude and are jointly significant over only a
narrow band near 600~m. That the transport-weighted limb
salinities agree to within a factor of 2.3 over 1993--2025 while the raw profiles
disagree by more than ten is an argument for weighting by the transport that
actually carries the salt rather than by depth.

ECCO is the counterexample, and against GLORYS12V1 it is a real one rather than
a lack of power. Its limb salinities are flat ($+0.005$~PSU per decade,
$p=0.55$, $N_{\text{eff}}=14.9$). On the matched 1993--2017 window its minimum
detectable $S_{\uparrow}$ trend is 0.020~PSU per decade, well below the
GLORYS12V1 rate of $+0.083$ over the same years, so ECCO could have resolved
that salinification and did not. Against ORAS5 the comparison is inconclusive:
the ORAS5 common-window rate, $+0.019$~PSU per decade, is itself below ECCO's
detection limit.

\subsection*{Is the widening contrast a water mass, or a redistribution?}

$\Delta S$ is transport-weighted, so it changes if the salinity field changes,
and also if transport is redistributed in depth at a completely fixed salinity
field. The two have to be separated before the change can be called a water-mass
signal. Recomputing $\Delta S$ with the salinity profile frozen at its record
mean, and again with the velocity profile frozen, does that at the full vertical
resolution of the section, and the answer differs between products.

In GLORYS12V1 the change is almost all salinity: the salinity-only trend is
$+0.068$~PSU per decade ($p=0.048$, $N_{\text{eff}}=4.8$) against a full trend
of $+0.072$, a share of 0.94, and the velocity-only trend is $+0.007$, not
significant and close to its own detection limit ($p=0.22$,
$N_{\text{eff}}=22.0$, detection limit 0.011).
In ORAS5 the partition is not resolved. The point estimates give a salinity
share of 0.66 and a velocity share of 0.39, so roughly a third of the widening
is a change in the vertical structure of the flow rather than in the water
masses. But the salinity-only series is so persistent (lag-1 autocorrelation
0.93) that its effective sample size falls to the floor of 3 imposed by the
significance test, and the smallest trend it could have resolved is
0.199~PSU per decade, eleven times the estimate itself. Its $p=0.45$ therefore
carries no information. The velocity-only term is the resolvable one
($+0.011$~PSU per decade, $p=0.0002$, $N_{\text{eff}}=31.9$, detection limit
0.005), so what can be said of ORAS5 is that a flow-structure contribution is
definitely present, not that the salinity contribution is absent. In GLORYS12V1
the salinity-only estimate also sits close to its own detection limit
($+0.068$ against 0.066, $N_{\text{eff}}=4.8$). In ECCO neither term is
resolvable ($-0.006$, detection limit 0.010; $+0.009$, detection limit 0.027)
and the full trend is indistinguishable from zero. The two shares sum
to slightly more than one because the split leaves a small covariance residual,
$-5$\% in ORAS5 and $-4$\% in GLORYS12V1.

The coarse component of that redistribution can be isolated further. The limbs
are defined by the sign of the section velocity, and that sign changes more than
once: besides the upper cell there is a northward abyssal limb of Antarctic
Bottom Water carrying 3.8~Sv in ORAS5, 4.8 in GLORYS12V1 and 2.2 in ECCO, that
is 20\%, 23\% and 15\% of the northward limb, and fresher than the upper cell by 0.21 to
0.49~PSU (Extended Data Fig.~2a,c). Writing $S_{\uparrow}=\sum_i w_i S_i$ over
the sub-limbs and holding either the weights or the sub-limb salinities at their
record means splits it. The partition term is small and of the opposite sign to
the total in both declining products: in ORAS5 the change in $S_{\uparrow}$ of
$+0.027$~PSU per decade is a water-mass term of $+0.030$ and a partition term of
$-0.003$, a share of 0.12; in GLORYS12V1, $+0.072$ splits into $+0.073$ and
$-0.002$, a share of 0.03. So the upper-to-abyssal redistribution is not what
drives the ORAS5 velocity term; that lives within the upper limb.

ECCO is the one product in which the coarse redistribution is significant and
active. Its upper-cell transport falls at $-0.60$~Sv per decade ($p=0.011$)
while its abyssal limb rises at $+0.36$ ($p=0.041$), and its upper-cell salinity
rises at $+0.022$~PSU per decade ($p=0.025$), the same sign as the other two
products. That last $p$ is basis-dependent: identifying the sub-limbs on
annual-mean profiles instead gives $+0.018$ and $p=0.18$, so the ECCO
cancellation is suggestive rather than established.
ECCO's flat whole-limb salinity is therefore the sum of a positive water-mass
term and a redistribution that cancels it, not an absence of change, and the
confound this section exists to test is demonstrably real even though it is not
what dominates elsewhere.

Excluding the abyssal cell altogether points the same way as the sub-limb split.
Truncating the section at the deep zero crossing and re-applying the throughflow
removal within the truncated layer, so the factorisation is again exact, gives
upper-cell $S_{\uparrow}$ trends of $+0.033$~PSU per decade in ORAS5
($p<10^{-4}$) and $+0.086$ in GLORYS12V1 ($p=0.033$), larger than the whole-limb
values, with transport shares of 0.112 and 0.019.

\subsection*{An observation-only test}

The claim that the northward limb is becoming saltier can be tested in a
salinity field that no ocean model integrates forward. We take EN4.2.2, an
objective analysis of profile observations, hold the velocity structure frozen
at each reanalysis record-mean profile, and let only the observed salinity vary
over 2005--2024 (Methods). Any trend is then a trend in observed water masses
under a fixed circulation. EN4 is on whole-degree latitudes, and we use the
34\degS{} row: the 35\degS{} row lies south of Cape Agulhas, where Africa no
longer closes the basin and a contiguous section runs into the Indian Ocean.

The limb asymmetry is reproduced, and more clearly than in the reanalyses'
own statistics. $S_{\uparrow}$ rises at $+0.026$~PSU per decade ($p=0.0005$)
under the ORAS5 velocity structure and $+0.028$ ($p=0.0003$) under the
GLORYS12V1 structure, while $S_{\downarrow}$ is flat in both ($+0.002$,
$p=0.09$), and the contrast widens at $+0.024$ to $+0.026$~PSU per decade
($p=0.0001$). The sign, the limb asymmetry and the significance therefore
survive in a salinity field that no assimilating model has touched.

The observed salinity field also supplies an observational estimate of $\Fovs$
itself. Combined with the frozen transport it gives $-0.080$~Sv under the ORAS5
velocity structure and $-0.128$~Sv under the GLORYS12V1 structure, both firmly
negative. This is not an independent measurement, since the transport is taken
from the reanalyses, but it shows the negative sign does not depend on the
reanalysis salinity fields. It also sharpens the disagreement about magnitude:
under ORAS5's own velocity structure the observed salinity field returns
$-0.080$~Sv against ORAS5's own $-0.038$~Sv, because EN4's limb contrast at this
section is 0.160~PSU against ORAS5's 0.068.

The rate is a different matter. Over the identical 2005--2024 window the
reanalyses give $S_{\uparrow}$ trends of $+0.056$ (ORAS5) and $+0.035$
(GLORYS12V1)~PSU per decade, against $+0.026$ and $+0.028$ from the
observations. ORAS5 runs about twice the observation-only rate while GLORYS12V1
is within 30\% of it. Part of that gap is expected by construction, since EN4
relaxes towards climatology where profiles are sparse and the core Argo array
does not sample below 2000~m, so the observation-only estimate is a lower bound.
The direction is robust; the ORAS5 rate should not be taken at face value.

\subsection*{The attribution does not depend on the choices we made}

\emph{Epoch selection.} Rather than quoting hand-picked epochs, we evaluate the
symmetric decomposition of equation~\eqref{eq:factorisation} for every
non-overlapping early and late epoch pair each record supports, at epoch lengths
of 10, 13 and 15 years: 2901 usable pairs in ORAS5 and 151 in GLORYS12V1. These
are not independent samples, and we treat them as a sensitivity analysis over
window choice rather than as evidence from a large sample. The median transport
share is 0.075 in ORAS5 and 0.038 in GLORYS12V1, the 95th percentile is 0.219
and 0.088, the maximum is 0.90 and 0.12, and the share is below 0.25 in 96.8\%
and 100\% of pairs (Fig.~\ref{fig:attribution}c). The transport term opposes the
total in 96.8\% and 94.7\% of pairs. The symmetric form uses the mean of the two
epochs as its coefficient, so it has no cross term and does not depend on which
epoch is the base point; the early-base and late-base conventions bracket it,
with median shares of 0.012 and 0.141 in ORAS5.

\emph{The definition of overturning strength.} Substituting the depth-space
streamfunction maximum for the limb transport, either integrated from the
surface or with the integration restarted at 250~m to exclude the wind-driven
surface cell, changes the mean overturning strength substantially, from 18.7 to
14.6 to 8.6~Sv in ORAS5, and raises the share to 0.111 [0.076, 0.151] and 0.183
[0.137, 0.240]; in GLORYS12V1 it becomes 0.014 and 0.040. The share therefore
rises monotonically as the definition of the overturning narrows, by a factor of
2.4 across the three definitions computed on this common basis (0.076, 0.111 and
0.183 for ORAS5), and the full range across every test in which the share is
identified, excluding the census tail, is 1\% to 18\% with individual intervals
reaching 27\%; the two largest values with a resolvable denominator are
0.18 over 1993--2025 in ORAS5 and 0.183 for the 250~m streamfunction definition
over its full record. It remains a
minority contribution throughout, but it is not definition-independent, and we
note that the factorisation we advocate gives a smaller share than the
construction it replaces.

\emph{The observing system.} The salinity field at 34.5\degS{} was barely
constrained before Argo, so we repeat everything from 2004 onwards. The
GLORYS12V1 decline survives ($-2.28$~mSv per year, $p=0.038$) and its transport
share rises, to 0.038 [0.004, 0.268], while remaining a minority. GLORYS12V1 has no documented stream join, but its record
crosses the Argo transition, and applying the same level-shift model at 2004
shows that most of its signal is a step rather than a trend: the residual trend
falls to $-1.90\pm1.07$~mSv per year ($p=0.098$, not significant) against a step
of $-0.061\pm0.022$~Sv ($p=0.014$), and the limb contrast behaves the same way
($+0.029\pm0.015$~PSU per decade, $p=0.071$, against a step of
$+0.105\pm0.031$, $p=0.004$). The attribution is unaffected, the step-adjusted
transport share being 0.045 and the step-adjusted fractional rate 23\% per
decade rather than 52\%, but the GLORYS12V1 \emph{decline} is a level shift
at the arrival of Argo and should not be read as a trend.

ORAS5 is also a splice, the consolidated
reanalysis to 2014 and the operational stream thereafter, and here too the test
is weaker than it looks. Fitting a trend together with a level shift at 2015 gives, over
the full record, a step in $\Fovs$ of $-0.017\pm0.011$~Sv ($p=0.15$); the
corresponding step in the limb salinity is larger relative to its own signal and
is significant, $+0.038\pm0.017$~PSU ($p=0.029$), and allowing for it reduces
the $S_{\uparrow}$ trend from $+0.025$ to $+0.020$~PSU per decade. Refitting over
2004--2025 alone gives a trend of $-1.36\pm1.53$~mSv per year that is not
significant, against $-2.46$ when the step is ignored. The ORAS5 Argo-era decline does not survive fitting the splice. Over the full record the step-adjusted trend is
$-1.17\pm0.21$~mSv per year ($p<10^{-4}$), essentially identical to the
consolidated-stream-only value of $-1.17$ ($p=5\times10^{-5}$), whose transport
share is 0.050 [0.024, 0.083]. So roughly 15\% of the raw ORAS5 full-record decline is
the splice, the decline survives correcting for it over the full record, and the
attribution is unaffected; but the 22-year Argo window alone cannot separate the
two.

\emph{The one window with neither confound.} 2004--2014 is inside the Argo era
and inside the ORAS5 consolidated stream at once. It is short, and it settles
nothing: over it the ORAS5 fractional-rate ratio is 0.152 with an $\Fovs$ trend that
is not resolvable, and GLORYS12V1 gives an $\Fovs$ trend of the opposite sign
($+7$\% per decade, $p=0.79$) against a transport falling at $-9$\% per decade
($p=0.0005$), so its fractional-rate ratio is 1.32; ECCO gives 0.30. Eleven years cannot
separate these quantities, and we report the window rather than omit it because
it is the only one that satisfies both homogeneity conditions.

\emph{The removal of the net transport.} Taking the net section transport out in
proportion to $|V|$, so that it is removed where the flow is rather than
uniformly over the area, changes $T$ by less than 0.6\% and leaves the
attribution intact: the transport share becomes 0.078 in ORAS5 and 0.012 in
GLORYS12V1, against 0.068 and 0.012 for the scheme used throughout.

\emph{The window.} Over 1993--2025, the longest window the two eddy-permitting
products share, both limb-salinity trends are resolvable ($+0.031$~PSU per
decade, $p=0.004$, in ORAS5 and $+0.071$, $p=0.017$, in GLORYS12V1, a factor of
2.3 apart rather than the factor of four over the shorter three-product window).
The ORAS5 transport share on that window is 0.180, and still a minority.

\emph{The section.} Recomputing the whole diagnostic at 25, 28, 30, 32 and
34.5\degS{} over the matched 1993--2025 window (Fig.~\ref{fig:sign}c), the
record mean is negative at every latitude in both products, and its interval
excludes zero everywhere except ORAS5 at 25\degS{} ($-0.004$~Sv). The magnitude
varies by a factor of 22 in ORAS5 across the full band and by 1.5 in
GLORYS12V1, and at 25\degS{} the two products disagree by a factor of 22, a
larger disagreement than any we report at 34.5\degS{}. (The two factors, 22.5
and 22.6, are close by coincidence.) Excluding 25\degS{}, the
ORAS5 range narrows to a factor of two. The negative sign is established in
ORAS5 only poleward of about 28\degS{}. We use 34.5\degS{} because that is
where the observing array sits and where the diagnostic is defined.

Two further quantities bound what the diagnostic means. The net volume
transport across the section is $-1.2$~Sv in ORAS5, $-1.2$ in GLORYS12V1 and
$-1.0$ in ECCO, and removing it in order to define an overturning component
shifts $\Fovs$ by $+0.009$, $+0.008$ and $+0.007$~Sv. Separately, the
overturning is not the only freshwater carrier at this section: the azonal
component is $+0.306$~Sv in ORAS5, $+0.304$ in GLORYS12V1 and $+0.155$ in ECCO,
opposite in sign to $\Fovs$ and, in the eddy-permitting products, several times
larger. Its trend is also opposite and, over each product's own record,
significant: $+1.84$~mSv per year in ORAS5 ($p=0.0002$, $N_{\text{eff}}=21.8$)
and $+6.68$ in GLORYS12V1 ($p=0.049$, $N_{\text{eff}}=8.5$, so marginal on few
effective degrees of freedom), both larger in magnitude than the $\Fovs$ trends
they accompany, against $+0.19$ ($p=0.95$) in ECCO, where the
detection limit is 14~mSv per year and the null is uninformative. The total freshwater transport
across 34.5\degS{} is therefore not moving the way $\Fovs$ alone implies, and a
reader assessing whether the basin is moving deeper into the self-reinforcing
regime must weigh both. $\Fovs$ is a stability diagnostic defined on the
overturning component; it is not the freshwater budget of the basin.

\subsection*{Implications for climate models}

CMIP6 historical simulations sit on the other side of zero, but not by a
majority that can be tested. Over 1950--1980, 16 of 25 models have a positive
mean $\Fovs$ at 34.5\degS{}, with a median of $+0.053$~Sv against a range of
$-0.205$ to $+0.405$~Sv (Fig.~\ref{fig:models}a). That count is not
distinguishable from a coin flip (binomial $p=0.23$), and taking one member per
model family leaves 18 independent models of which 10 are positive
($p=0.81$), or 9 if the other member of the one internally inconsistent family
is chosen ($p=1.00$). The defensible statement is about location rather than
majority: the model median is $+0.053$~Sv while all three recomputed
observational estimates lie between $-0.038$ and $-0.132$~Sv, below 72\%, 80\%
and 88\% of the models respectively. This confirms a bias documented
previously\cite{Mecking2017,vanWestenDijkstra2024,Weijer2019} and benchmarks it
against three consistently computed estimates; it does not discover it. The
model means are over 1950--1980 while the reanalyses cover modern decades, and
since ORAS5 declines by 1.37~mSv per year the two windows are not equivalent;
the comparison should be read as a statement about a persistent model offset,
not as a like-for-like difference of contemporaneous states.

Whether that bias matters for projected rates is a separate question, and the
test that has been used to answer it has no power. Correlating model mean
$\Fovs$ against a 19-year overturning trend at 26.5\degN{} over 2005--2023
gives $r=-0.04$ ($p=0.88$) across 16 models, but the smallest correlation that
sample could reach significance with is $|r|=0.50$, and 19-year overturning
trends within a single model's large ensemble already span $-3.5$ to
$+2.4$~Sv per decade across 50 members of MPI-ESM1-2-LR. A null from such a test
is guaranteed by internal variability whatever the forced relationship. The
counting statistic behaves the same way: 2 of 6 models with negative $\Fovs$
weaken faster than GLORYS12V1 (whose own rate at 26.5\degN{} is $-0.70$~Sv per
decade) against 6 of 10 with positive $\Fovs$, a Fisher exact $p$ of 0.61, with the same counts against the RAPID rate.

Repeating the test on the forced response, where the signal is large, remains
negative but is now informative. Across 21 models with historical plus ssp585
projections to 2100, mean $\Fovs$ and the overturning change from 1850--1900 to
2081--2100 give $r=-0.20$ ($p=0.40$, 95\% CI $-0.58$ to $+0.26$) against a
smallest reachable $|r|$ of 0.43 (Fig.~\ref{fig:models}b); on fractional rather
than absolute change $r=+0.30$ ($p=0.18$). Leaving out one model at a time
moves the first between $-0.41$ and $-0.05$ and the second between $+0.06$ and
$+0.45$, so neither is stable.

The reason the two have opposite signs is itself the one significant result in
this comparison. Model mean $\Fovs$ and the pre-industrial overturning strength
are strongly correlated ($r=+0.73$, $p=1.5\times10^{-4}$, 95\% CI $+0.44$ to
$+0.89$, $n=21$): models with a stronger overturning carry a more positive
$\Fovs$, which is what equation~\eqref{eq:factorisation} predicts given that
these models mostly carry $\Delta S<0$. Absolute overturning change is
mechanically tied to the base state while fractional change is not, which is
why the two rate tests disagree in sign. It also suggests that the CMIP6
$\Fovs$ bias is in part a symptom of the overturning strength bias rather than
an independent error in the freshwater budget.

The observed rate is itself weakly constrained. Over 2005--2023 the RAPID array
gives $-0.89$~Sv per decade at 26.5\degN{}, with an unadjusted $p$ of 0.093
that becomes 0.21 once the effective sample size is accounted for
($N_{\text{eff}}=11.9$). The direct observations do not yet establish a
significant decline over the era in which the reanalyses are best
constrained\cite{Jackson2024,Terhaar2025}.

\clearpage

\begin{figure}[t]
  \centering
  \includegraphics[width=\textwidth]{Fig1_sign.pdf}
  \caption{\textbf{The sign of the overturning freshwater transport at
  34.5\degS{}.}
  \textbf{a}, Annual-mean $\Fovs$. ORAS5, GLORYS12V1 and ECCO-V4r4 are computed
  here from monthly fields on the contiguous Atlantic section; SODA 3.15.2 is
  carried over from an annual-field pipeline, is not on the same basis, and is
  not used in any quantitative claim.
  \textbf{b}, Record mean (circles) and most recent decade (diamonds) with 95\%
  circular block-bootstrap intervals, at a 5-year block for the record mean and
  a 3-year block for the decade. The pale bar spans the intervals obtained for
  block lengths of 2 to 20 years. Between-product spread is 0.094~Sv, wider than
  the widest of these intervals (0.073~Sv, GLORYS12V1 across block lengths) and
  1.5 times the 5-year intervals quoted in the text.
  \textbf{c}, Mean $\Fovs$ over 1993--2025 against section latitude, with the
  same intervals. The open marker is the one case whose interval includes zero.}
  \label{fig:sign}
\end{figure}

\begin{figure}[t]
  \centering
  \includegraphics[width=\textwidth]{Fig2_attribution.pdf}
  \caption{\textbf{The decline is carried by the salinity contrast, not by the
  transport.}
  \textbf{a}, The two factors of equation~\eqref{eq:factorisation}, each as a
  percentage anomaly from its own record mean, so that a factor which barely
  moves can be compared with one that moves a great deal. The ORAS5 contrast
  passes close to zero early in the record, which is why its percentage anomaly
  is large there.
  \textbf{b}, Trend in $\Fovs$ split into transport and salinity contributions
  with 95\% block-bootstrap intervals, and the covariance residual, over each
  product's own record and over the Argo era.
  \textbf{c}, Transport share of the change over every non-overlapping epoch
  pair the record supports (2901 for ORAS5, 151 for GLORYS12V1) at epoch lengths
  of 10, 13 and 15 years, treated as a sensitivity analysis over window choice.
  Ticks mark the median, 90th and 95th percentiles and the maximum; the
  percentage is the fraction of pairs below the dashed 0.25 line.}
  \label{fig:attribution}
\end{figure}

\begin{figure}[t]
  \centering
  \includegraphics[width=\textwidth]{Fig3_mechanism.pdf}
  \caption{\textbf{The northward limb salinifies; the southward limb does not.}
  \textbf{a}, Transport-weighted salinity of the northward limb (solid) and the
  southward limb (dashed), as anomalies from their record means.
  \textbf{b}, Zonal-mean salinity trend against depth at the section over
  1993--2025, drawn bold where the pointwise Santer-adjusted $p$ is below 0.05.
  The panel is truncated at 3000~m; the abyssal limb discussed in the text
  begins below 4000~m.
  Adjacent levels are strongly correlated, so no field-significance claim is
  made from this panel.
  \textbf{c}, Transport share of the $\Fovs$ trend under every robustness test,
  with 95\% block-bootstrap intervals: the full record, the Argo era, two
  alternative definitions of the overturning strength, and, for ORAS5, the
  homogeneous consolidated stream alone.}
  \label{fig:mechanism}
\end{figure}

\begin{figure}[t]
  \centering
  \includegraphics[width=\textwidth]{Fig4_models.pdf}
  \caption{\textbf{CMIP6 lies above the observational estimates, but cannot be
  sorted by the sign of $\Fovs$.}
  \textbf{a}, Mean $\Fovs$ at 34.5\degS{} over 1950--1980 in 25 CMIP6 models,
  with the reanalysis estimates marked; SODA is shown for continuity with
  Fig.~\ref{fig:sign} and is not used quantitatively. No offset or bias
  correction is applied to either models or reanalyses.
  \textbf{b}, Model mean $\Fovs$ against the overturning change at 26.5\degN{}
  from 1850--1900 to 2081--2100 in the 21 models with historical plus ssp585
  projections ($r=-0.20$, $p=0.40$, smallest reachable $|r|=0.43$). This
  replaces a 19-year single-member trend test, whose null result is guaranteed
  by internal variability: 19-year trends span $-3.5$ to $+2.4$~Sv per decade
  across 50 members of the MPI-ESM1-2-LR large ensemble.}
  \label{fig:models}
\end{figure}
